% Pursuing Stacks (cotes 134-2 to 134-8): archive cote 134-6, site page 49 (PDF page 50), typescript page 367.
% typescript. Draft compared: ps5.tex:3785-3857.
% AI-transcribed by Claude Opus 5.5 (claude-opus-5-5, adaptive thinking, effort high), Message Batches API, batch ps3 (msgbatch_01XdFmuVVBAcgc7EFxdtCqsj), 2026-11-12; reviewed against both supplied scan halves; see CHECK_DECISIONS.tsv.
\scanpage{134-6/49}
\mspage{367}

17.6.

100.\quad Yesterday I introduced the notion of an \uline{integrator} for any small
category B, to be just a projective resolution of $\underline{Z}_B$ in the category $B^{\hand{\wedge}}_{ab}$
of all abelian presheaves on B, where $\underline{Z}_B$ denotes the constant presheaf
with value $\underline{Z}$ . Such an object in Ch.$(B^{\hand{\wedge}}_{ab})$ exists, due to the existence
of sufficiently many projectives in $B^{\hand{\wedge}}_{ab}$ , and it is unique up to homoto-%
pism of \inserted[typed]{augmented} chain complexes, which encourages us to denote it by a canonizing
symbol, namely
\[ (1) \qquad L^{B}_{\bullet} \hand{\to} \underline{Z}_B \quad . \]
As will become cle\over[typed]{q}{a}r in the sequel, $L^{B}_{\bullet}$ can be viewed as embodying \uline{homolo-%
gy properties} of B, i.e. of the topos associated to B (whose category of
sheaves of sets is $B^{\hand{\wedge}}$). The way we hit upon it though was in order to ob-%
tain a ``computational'' way for computing \uline{cohomology} of B (i.e. of the asso-%
ciated topos) with coefficients in any abelian \xout{sheaf} presheaf F in $B^{\hand{\wedge}}_{ab}$ ,
by a\xout{n} canonical isomorphism
\[ (2) \qquad R\Gamma_B(F) \overset{\hand{\sim}}{\hand{\longrightarrow}} \mathrm{Hom}^{\bullet}(L^{B}_{\bullet}, F) \]
in the derived category $D^{\bullet}(Ab)$, where $\mathrm{Hom}^{\bullet}$ denotes the cochain complex
ob\over[typed]{y}{t}ained by applying Hom componentwise. Passing to the cohomology groups
of both members, this gives rise to
\[ (3) \qquad H^{i}(B,F) \hand{\simeq} H^{i}\mathrm{Hom}^{\bullet}(L^{B}_{\bullet} , F) \quad . \]
T\tikz[baseline=(hm.base)]{\node[inner sep=0pt] (hm) {h};\draw ([xshift=-.4ex,yshift=.45ex]hm.north)--([xshift=.4ex,yshift=.45ex]hm.north);\draw ([xshift=-.25ex,yshift=.45ex]hm.north)--([xshift=-.25ex,yshift=.75ex]hm.north);\draw ([xshift=.25ex,yshift=.45ex]hm.north)--([xshift=.25ex,yshift=.75ex]hm.north);\draw ([xshift=-.3ex,yshift=-.3ex]hm.south)--([xshift=-.3ex,yshift=-.8ex]hm.south);\draw ([xshift=-.3ex,yshift=-.8ex]hm.south) .. controls ([xshift=.15ex,yshift=-.4ex]hm.south) and ([xshift=.4ex,yshift=-.4ex]hm.south) .. ([xshift=.4ex,yshift=-.8ex]hm.south);}e \over[typed]{d}{d}esignation ``computational'' takes a rather concrete meaning, when we
choose $L^{B}_{\bullet}$ to have it's components in the infinitely additive envelope
Addinf(B) of B, which (as w\over[typed]{i}{e} saw yesterday) can be viewed as a full subca-%
tegory of $B^{\hand{\wedge}}_{ab}$ , made up with projectives, and such that any object in
$B^{\hand{\wedge}}_{ab}$ is quotient of an object coming from Addinf(B); \xout{which} \inserted[typed]{\uline{this}} insures that
there exist indeed integrators \over{w}{w}\over{h}{h}ich are \uline{``quasi-special''} i.e. which are
made up with objects of Addinf(B), and hence can be interpreted as chain
complexes of this additive category. Thus any component $L_n$ can now be
written, in an essentially canonical way, as
\[ (4) \qquad L_n = \hand{\bigoplus}_{i\hand{\in} I_n} \underline{Z}^{(b_i)} \quad , \]
Where
\[ (5) \qquad (b_i)_{i\hand{\in} I_n} \]
is a family of objects of B indexed by $I_n$ (NB for simplicity of notations,
we assume the $I_n$'\hand{s}\deleted{s} mutually disjoint, otherwise we shou\over[typed]{b}{l}d write the
\margin{49}

% Visual collation: both supplied image halves inspected; see CHECK_DECISIONS_134_7_8.tsv.
